%============================================================================
%  FIRST DRAFT — Prioritized memory transfer: a predictive-coding primitive
%  for continual learning.  (Saighi & Rozenberg)
%
%  Format follows the PLOS-style example layout (../example layout/manuscript.tex);
%  narrative ordering follows Tang et al. 2023 (results-forward, theorems stated
%  inline in a Model/analysis block, full proofs in Materials and methods).
%  Full derivations: ../maths/detailed_proofs.md ; house identities: ../maths/cheatsheet.md.
%
%  Build:  pdflatex draft && bibtex draft && pdflatex draft && pdflatex draft
%============================================================================
\documentclass[10pt,letterpaper]{article}
\usepackage[top=0.85in,left=2.75in,footskip=0.75in]{geometry}

\usepackage{amsmath,amssymb,amsthm}
\usepackage{graphicx}
\usepackage{algorithm}
\usepackage{algpseudocode}
\usepackage{changepage}
\usepackage{subcaption}
\usepackage{textcomp,marvosym}
\usepackage{cite}
\usepackage{nameref,hyperref}
\usepackage[right]{lineno}
\usepackage[nopatch=eqnum]{microtype}
\DisableLigatures[f]{encoding = *, family = * }
\usepackage[table]{xcolor}
\usepackage{array}

% ---- theorem environments (the paper's formal results) ----
\theoremstyle{plain}
\newtheorem{theorem}{Theorem}
\newtheorem{lemma}{Lemma}
\newtheorem{proposition}{Proposition}
\newtheorem{corollary}{Corollary}
\theoremstyle{remark}
\newtheorem{remark}{Remark}

% ---- text layout (PLOS-like) ----
\raggedright
\setlength{\parindent}{0.5cm}
\textwidth 5.25in
\textheight 8.75in

\usepackage[aboveskip=1pt,labelfont=bf,labelsep=period,justification=raggedright,singlelinecheck=off]{caption}
\renewcommand{\figurename}{Fig}

\bibliographystyle{bibstyle}
\makeatletter
\renewcommand{\@biblabel}[1]{\quad#1.}
\makeatother

% ---- header / footer ----
\usepackage{fancyhdr}
\pagestyle{fancy}
\fancyhf{}
\rfoot{\thepage}
\renewcommand{\headrulewidth}{0pt}
\renewcommand{\footrule}{\hrule height 2pt \vspace{2mm}}
\lfoot{Draft \today}

% ---- lightweight figure placeholder (no image files needed yet) ----
\newcommand{\figplaceholder}[1]{%
  \begin{center}\setlength{\fboxsep}{14pt}%
  \fbox{\parbox{0.82\textwidth}{\centering\itshape [figure placeholder]\\[2pt] #1}}%
  \end{center}}

% ---- notation macros (consistent with cheatsheet.md, house identities H1--H12) ----
\newcommand{\R}{\mathbb{R}}
\newcommand{\eps}{\varepsilon}
\newcommand{\xT}{x_T}
\newcommand{\xS}{x_S}
\newcommand{\WT}{W_T}
\newcommand{\WS}{W_S}
\newcommand{\MT}{M_T}
\newcommand{\MS}{M_S}
\newcommand{\ST}{S_T}
\renewcommand{\SS}{S_S}
\newcommand{\NS}{N_S}
\newcommand{\FS}{F_S}
\newcommand{\FSeq}{F_S^{\mathrm{eq}}}
\newcommand{\FT}{F_T}
\newcommand{\epsT}{\eps_T}
\newcommand{\epsS}{\eps_S}
\newcommand{\epsTS}{\eps_{TS}}
\newcommand{\piT}{\pi_T}
\newcommand{\piS}{\pi_S}
\newcommand{\piTS}{\pi_{TS}}
\newcommand{\piST}{\pi_{ST}}
\newcommand{\tauT}{\tau_T}
\newcommand{\tauS}{\tau_S}
\newcommand{\UT}{U_T}
\newcommand{\tr}{^{\!\top}}

\begin{document}

\vspace*{0.2in}
\begin{flushleft}
{\Large
\textbf\newline{Prioritized memory transfer: a predictive-coding primitive for continual learning}
}
\newline
\\
Saighi Paul*\textsuperscript{1},
Rozenberg Marcelo\textsuperscript{1,2}
\\
\bigskip
\textbf{1} Universit\'e Paris-Saclay, CNRS, Laboratoire de Physique des Solides, 91405, Orsay, France
\\
\textbf{2} Universit\'e Paris-Cit\'e, CNRS, Integrative Neuroscience and Cognition Center, 75006, Paris, France
\\
\bigskip
* paul.saighi@universite-paris-saclay.fr
\end{flushleft}

%============================================================================
\section*{Abstract}
%============================================================================
Biological memory is consolidated during sleep, when the brain replays stored experience with no
external teacher and, remarkably, appears to favour what has not yet been consolidated. Reproducing
this in artificial networks is the crux of continual learning: how can a network rehearse its own
memories, and how can that rehearsal be \emph{prioritized} toward what a downstream network still
lacks, without an external scheduler holding the list of what to replay? We show that predictive
coding answers both questions at once. Two linear associative memories are coupled in a hierarchy: a
frozen \emph{teacher} above an empty, plastic \emph{student}. Flipping the sign of a single
interface precision turns ordinary recall into replay. Under this ``reversed precision'' the student's
own prediction error drives the teacher into exactly the memories the student cannot yet predict; the
student learns them; the drive there vanishes; the process stops itself. We prove that the drive is
the gradient of the student's own surprise, so priority is not imposed but derived --- the student's
error \emph{is} the priority signal --- and that at a precise matching condition the coupled dynamics
form an exact free-energy saddle. Because the substrate is linear, transfer acquires a memory
\emph{subspace} rather than named patterns. Transfer is a primitive that composes: interleaving two
teachers merges their subspaces and cancels the crosstalk between correlated memories, and a
buffer\,$\to$\,consolidate\,$\to$\,storage loop built on it retains a streamed sequence of correlated
memories where a buffer-only control catastrophically forgets. The result is a compact, analytically
transparent, biologically suggestive account of prioritized replay and its use for continual learning.

%============================================================================
\section*{Author summary}
%============================================================================
When we sleep, the brain spontaneously replays recent experience, and this replay helps move memories
into long-term storage without erasing what is already there. Artificial neural networks are notoriously
bad at this: teaching them something new tends to overwrite the old, a failure called catastrophic
forgetting. Rehearsing old memories fixes it, but only if the network can somehow retrieve those
memories by itself and, ideally, spend its effort on the memories that still need consolidating. We
describe a simple two-network model in which this happens automatically. A frozen ``teacher'' network
holds memories; an empty ``student'' network learns them. By reversing the sign of a single connection
parameter, the pair switches from ordinary recall into a replay mode in which the teacher is pushed
toward whatever the student has not yet learned --- and the push is generated entirely by the student's
own prediction error, so no outside controller is needed to decide what to replay. We prove that this
push always climbs the student's ``surprise'' and switches itself off once a memory is learned. Stacking
this transfer step into a small loop lets a network absorb a stream of new, overlapping memories while
keeping the earlier ones, offering a biologically plausible route around catastrophic forgetting.

\linenumbers
\newpage

%============================================================================
\section{Introduction}
%============================================================================
A system learns continually when it can absorb new information over time without discarding what it
already knows. Artificial neural networks fail at this in a characteristic way: training on a new task
or pattern in isolation rapidly erodes the representations built for earlier ones, the phenomenon of
\emph{catastrophic forgetting} \cite{mccloskey1989catastrophic,robins1995catastrophic,kirkpatrick2017ewc}.
In associative memories the failure is concrete --- storing a new activity pattern sequentially degrades
the previously stored ones.

The oldest and most robust remedy is \emph{rehearsal}: periodically retrain the network on its earlier
memories so their representations are reinforced while new ones are added
\cite{robins1995catastrophic}. Rehearsal has a striking biological counterpart. During sleep and quiet
waking the hippocampus spontaneously reactivates recent experience, and this replay is causally
implicated in memory consolidation \cite{wilson1994reactivation,diekelmann2010memory}. The
complementary-learning-systems account casts a fast hippocampal store and a slow neocortical store in
exactly the teacher/apprentice roles that rehearsal requires: the fast system replays to train the slow
one \cite{mcclelland1995cls,kumaran2016what}. Two questions, however, are usually left open. First,
rehearsal presupposes that the network can \emph{retrieve} its own memories to replay them, rather than
reading them from an external list. Second, replay is not uniform --- consolidation should concentrate
on what the slow system has not yet acquired. \emph{How is autonomous replay generated, and how is it
prioritized, without an external scheduler?}

In earlier work we addressed the first question within a single continuous Hopfield network: plastic
self-inhibition lets the network shallow out already-visited attractors and so recover its full set of
stored patterns from its own dynamics, which is enough to drive a continual-learning loop
\cite{saighi2025autonomous}. Here we address the second, and we do so by moving from one network to two
coupled by \emph{predictive coding}. Predictive coding models a cortical hierarchy as levels that each
try to predict the level below and pass up only the residual error, learning by descending a
variational free energy through purely local computations
\cite{raoballard1999,whittington2017,millidge2022review}; Tang, Barron and Bogacz recently showed that a
recurrent predictive-coding network with covariance learning is a faithful, biologically grounded
associative memory \cite{tang2023covpcn}, and that substrate is the one we build on. Coupling two such
memories exposes a quantity a single network does not have: the prediction error of one network about
the other. That error is precisely a measure of what one network cannot yet predict, and we use it as
the priority signal --- an idea that mirrors prioritized experience replay in reinforcement learning,
where transitions are replayed in proportion to their error \cite{schaul2016prioritized}, but which here
falls out of the network's own free energy rather than being engineered. The free-energy reading also
connects the mechanism to active inference, in which agents act to minimize (variational) free energy
\cite{friston2010free,friston2017process}; we will draw that connection precisely where it is exact and
flag it as merely structural where it is not \cite{millidge2021whence}.

Our contribution is one mechanism presented as a ladder of three rungs, each built on the one below.

\begin{enumerate}
\item \textbf{Transfer.} One teacher $T$ drives one student $S$. A single signed interface precision
selects the phase: positive is wake/recall, negative is sleep/replay. In replay the student's
prediction error steers the teacher into the memories the student lacks; the student learns them; the
drive self-extinguishes. We prove the drive is gradient ascent on the student's own surprise, and that
the coupled flow is an exact free-energy saddle at a precise precision-matching condition.
\item \textbf{Union of subspaces.} Because transfer is a reusable primitive, coupling a student to two
teachers in alternation merges their memory subspaces. Interleaving one teacher per replay bout cancels
the crosstalk that correlation would otherwise create between the two sets of memories.
\item \textbf{Continual learning.} A small buffer\,$\to$\,consolidate\,$\to$\,storage loop, with the
consolidation step performed by interleaved transfer, retains a stream of correlated memories where a
buffer-only control forgets all but the most recent.
\end{enumerate}

Throughout, the substrate is a \emph{linear} covariance predictive-coding memory. Linearity is a
deliberate simplification: it makes the whole mechanism analytically transparent, at the cost that what
transfers is a memory \emph{subspace} (an effective rank) rather than a set of named episodes. We state
this scope plainly and return to it in the Discussion.

%============================================================================
\section{Model and analysis}
%============================================================================
We first fix the model and then state its formal properties. Following Tang et al.\ \cite{tang2023covpcn}
and the predictive-coding convention, results are stated here and proved in full in Materials and methods.

\subsection{A two-population predictive-coding memory}
\label{sec:model}
Both networks live in the same space $\R^{d}$, one coordinate per neuron. Each is a linear
predictive-coding associative memory: a population with recurrent weights $W$ (zero diagonal, so no
neuron predicts itself) tries to reconstruct its own state, producing the \emph{mismatch}
\begin{equation}
\label{eq:mismatch}
M = I - W, \qquad \eps = M x, \qquad S = M\tr M,
\end{equation}
where $\eps$ is the reconstruction error and $S$, its Gram matrix, is the symmetric positive
semi-definite \emph{self-error operator}. A pattern $m$ is stored when it is a fixed point of the
reconstruction, $M m = 0$; the stored memories therefore span the kernel of $M$, a flat linear
manifold (house identity H1--H2 of the accompanying notes).

We couple a \textbf{teacher} $T$ and a \textbf{student} $S$ (Fig~\ref{fig:schematic}). The teacher sits
higher in the hierarchy: it is a fully trained memory, its weights $\WT$ frozen, and it is confined to
its own memory manifold. The student sits below: it starts empty ($\WS=0$) and is plastic. Three errors
drive the pair --- the teacher's self-error $\epsT=\MT\xT$, the student's self-error $\epsS=\MS\xS$, and
the \emph{interface error}
\begin{equation}
\label{eq:interface}
\epsTS = \xS - \xT,
\end{equation}
the disagreement between the two states, which the identity interface reports at the student's level.
Each error is weighted by a precision (an inverse variance): $\piT,\piS>0$ are the self precisions,
$\piTS>0$ is the precision with which the student trusts its input from the teacher, and $\piST$ is the
teacher's interface precision --- the one quantity we allow to be \emph{signed}. The state and weight
dynamics are
\begin{align}
\tauT\,\dot\xT &= -\,\piT \MT\tr\epsT \;+\; \piST\,\epsTS \;+\; \xi, \label{eq:teacher}\\
\tauS\,\dot\xS &= -\,\piTS\,\epsTS \;-\; \piS \MS\tr\epsS, \label{eq:student}\\
\dot\WS &= \eta\,\piS\,\epsS\,\xS\tr, \qquad \mathrm{diag}(\WS)\equiv 0, \label{eq:learn}
\end{align}
with $\xi$ a small isotropic noise on the teacher and the timescales ordered
$\tauS \ll \tauT \ll 1/\eta$: the student perceives fast, the teacher drifts slowly, the weights change
slowest of all. The teacher is additionally held to a fixed norm, so it can reorient but not run away.

\begin{figure}[h!]
\centering
\figplaceholder{\textbf{Two coupled predictive-coding memories.} Frozen teacher $T$ (above) over an
empty, plastic student $S$ (below); each has recurrent weights $W$ and a self-error population
($\epsT$, $\epsS$), and the identity interface reports the disagreement $\epsTS=\xS-\xT$. The teacher's
interface precision $\piST$ is signed: $\piST>0$ (wake/recall) pulls both toward agreement --- no
transfer; $\piST<0$ (sleep/replay) reverses the teacher's drive, pushing it toward whatever the student
cannot predict. One sign is the whole phase.}
\caption{\textbf{Model architecture and the wake/sleep sign flip.} See Section~\ref{sec:model}.}
\label{fig:schematic}
\end{figure}

\paragraph{Reversed precision: one knob for wake and sleep.} Everything hinges on the sign of $\piST$
in Eq.~\eqref{eq:teacher}. When $\piST>0$ the teacher, like the student, is pulled to \emph{reduce} the
interface error: both networks seek agreement, the student completes the teacher's pattern, and nothing
is transferred. This is ordinary recall --- \emph{wake}. When $\piST<0$ the same term becomes a
\emph{reversed precision}: the teacher is pushed to \emph{increase} the disagreement, a drive to seek
out states the student cannot match. This is \emph{sleep}, or replay. No separate gate switches the
phase; the sign of a single precision \emph{is} the phase. This is the paper's pivot, and the rest of
the analysis simply works out its consequences.

\paragraph{Surprise and free energy.} Each network carries a free energy, the precision-weighted sum of
its squared errors,
\begin{equation}
\label{eq:freeenergy}
\FS = \tfrac{\piTS}{2}\,\|\epsTS\|^{2} + \tfrac{\piS}{2}\,\|\epsS\|^{2},
\qquad
\FT = \tfrac{\piT}{2}\,\|\epsT\|^{2},
\end{equation}
the network's \emph{surprise} at its current state. We keep two words strictly apart throughout, and the
reader will save themselves confusion by doing the same:

\begin{remark}[Novelty versus surprise]
\label{rem:terminology}
\emph{Novelty} is a property of the student's \emph{weights}: which directions it cannot yet predict.
It is carried by an operator, $\NS$ below, and its eigenvalues $n_k$. \emph{Surprise} is a property of a
\emph{state}: how badly the current configuration is predicted, measured by the scalar free energies
$\FS,\FT$. Novelty is the spectrum; surprise is the score. There is no ``teacher's surprise'' other than
$\FT$.
\end{remark}

Two inequalities delimit the regime we operate in, and we state them as such rather than as hypotheses
of the theorems: the input guard $\piTS>\piS$ (the student trusts its teacher input over its own
recurrent prior, which prevents it from confabulating), and a stability guard on the \emph{magnitude} of
the reversed precision, made precise in Lemma~\ref{lem:normal}.

\subsection{Inference and learning descend a free energy}
The first property justifies calling $\FS$ a free energy at all: the student's dynamics are gradient
descent on it.

\begin{proposition}[Free-energy descent]
\label{prop:descent}
The student's state relaxation \eqref{eq:student} is exact gradient descent on $\FS$,
$\tauS\dot\xS=-\nabla_{\xS}\FS$, and its zero-diagonal weight update \eqref{eq:learn} is projected
gradient descent on $\FS$; along both flows $\FS$ is non-increasing.
\end{proposition}

The state gradient is $\nabla_{\xS}\FS=\piTS\epsTS+\piS\MS\tr\epsS$, which is exactly minus the
right-hand side of Eq.~\eqref{eq:student}; the weight gradient is
$\nabla_{\WS}\FS=-\piS\epsS\xS\tr$, and removing its diagonal is an orthogonal projection, so descent
survives it (H9; proof in Materials and methods, verified numerically to machine precision).

\subsection{The novelty operator}
Because the student is fast (\,$\tauS\ll\tauT$\,), on the teacher's timescale it always sits at its own
equilibrium. Solving $\nabla_{\xS}\FS=0$ for the settled student gives the operator that governs
everything downstream.

\begin{lemma}[Novelty operator]
\label{lem:novelty}
At the fast-student equilibrium the student's state is
$\xS^{*}=(I-\NS)\,\xT$ and the interface error is $\epsTS=-\NS\xT$, with
\begin{equation}
\label{eq:novelty}
\NS \;=\; \piS \SS\,(\piTS I + \piS \SS)^{-1}.
\end{equation}
$\NS$ is symmetric positive semi-definite, shares the eigenbasis of $\SS$, and annihilates every
learned direction: if $\MS u=0$ then $\NS u=0$.
\end{lemma}

The reading is clean. The student copies back the part of the teacher's state it can predict,
$\xS^{*}=(I-\NS)\xT$, and leaves the part it cannot, so the residual disagreement $\epsTS=-\NS\xT$ is
\emph{exactly the novel part of the teacher's state} (H5--H6). In the shared eigenbasis
$\NS u_k=n_k u_k$ with novelty eigenvalues
\begin{equation}
\label{eq:spectrum}
n(\mu)=\frac{\piS\,\mu}{\piTS+\piS\,\mu}\in[0,1),
\end{equation}
where $\mu\ge 0$ is the corresponding eigenvalue of $\SS$: a fully learned direction ($\mu=0$) has
$n=0$, and novelty grows monotonically with how badly the student mispredicts the direction. $\NS$ is a
\emph{novelty meter} built from the student's weights alone.

\subsection{The surprise identity}
Novelty (a weight-level spectrum) and surprise (a state-level scalar) are not two mechanisms but two
readings of one quantity. This is the bridge, and it is exact.

\begin{corollary}[Surprise identity]
\label{cor:surprise}
At the fast-student equilibrium the student's surprise about the teacher's state is the teacher's state
scored by the novelty meter:
\begin{equation}
\label{eq:surpriseid}
\FSeq(\xT)\;\equiv\;\min_{\xS}\FS(\xS,\xT)\;=\;\tfrac{\piTS}{2}\,\xT\tr\NS\,\xT,
\qquad\text{hence}\qquad
\nabla_{\xT}\FSeq=\piTS\,\NS\,\xT.
\end{equation}
\end{corollary}

The proof is a two-line computation --- per eigendirection the interface term contributes $n^{2}$ and
the self term $n(1-n)$, which sum to $n$ (H8) --- or a one-line application of the envelope theorem,
since $\xS^{*}$ minimizes $\FS$. Either way, the settled student's free energy is literally
$\sum_k n_k c_k^{2}$ with $c_k=u_k\tr\xT$: novelty priced out on where the teacher's state actually
sits. The gradient we will need in Theorem~\ref{thm:prioritized} follows immediately.

\subsection{The student steers the teacher up the student's surprise}
We can now read the teacher's sleep dynamics. Substituting the settled interface error
$\epsTS=-\NS\xT$ into Eq.~\eqref{eq:teacher} and writing $\piST=-|\piST|$ for the reversed precision,
\begin{equation}
\label{eq:teacher_reduced}
\tauT\,\dot\xT \;=\; \underbrace{-\,\piT \ST\,\xT}_{\text{holds its own memory}}
\;+\; \underbrace{|\piST|\,\NS\,\xT}_{\text{the student's push } u},
\end{equation}
the teacher's own self-pull plus a push $u=|\piST|\NS\xT$ produced entirely by the student. This push
is the heart of the paper.

\begin{theorem}[Prioritized transfer]
\label{thm:prioritized}
The student's push $u=|\piST|\,\NS\,\xT$ has three equivalent readings.
\emph{(i) Per direction:} in the student's eigenbasis each component of the teacher's state grows at a
rate set by that direction's novelty, $\dot c_k=(|\piST|/\tauT)\,n_k\,c_k$; a learned direction
($n_k=0$) is frozen, a novel one is amplified, and the more novel, the faster.
\emph{(ii) Steepest ascent:} the push points along the gradient of the student's surprise,
$u=(|\piST|/\piTS)\,\nabla\FSeq$, hence it is the direction of fastest increase of $\FSeq$.
\emph{(iii) Self-limiting:} as the student learns a direction ($\|\MS u_k\|\!\to\!0$, by
Proposition~\ref{prop:descent}) its novelty $n_k\!\to\!0$, so the push there fades to zero.
Consequently the full teacher flow is
$\tauT\dot\xT=-\nabla\FT+(|\piST|/\piTS)\,\nabla\FSeq$: the teacher descends its own surprise while
ascending the student's.
\end{theorem}

In words: the teacher's state is passed through the student's novelty meter, and the meter amplifies
whatever the student finds novel and ignores whatever it has learned. Reading (ii) says this is not an
arbitrary drift but the steepest local climb of the student's own surprise --- priority is
\emph{derived}, not imposed. Reading (iii) says the climb flattens the very landscape it climbs: every
direction the teacher is dragged into gets learned, its hill sinking to a plain, so the process stops
itself when the student has acquired everything the teacher holds. Two honest caveats travel with the
theorem: the ascent is local (steepest from where the teacher currently is, not a beeline to the globally
most novel direction), and a direction the teacher does not currently occupy ($c_k=0$) receives no push
--- it is the noise $\xi$ that seeds those, and the reversed-precision drive that then amplifies them.

\subsection{Staying on the memory manifold}
The push must not tip over into chasing noise off the manifold. The teacher's self-pull prevents this,
provided the reversed precision is not too strong.

\begin{lemma}[Off-manifold stability]
\label{lem:normal}
Under the conservative guard $|\piST| < \piT\,\sigma_{\min}^{2}$, where $\sigma_{\min}^{2}$ is the
smallest nonzero eigenvalue of $\ST$, the block of the dynamics normal to the teacher's memory manifold
is strictly contracting, so the teacher consolidates on-manifold and stays quiescent. Above the guard
it is driven off-manifold into noise. (The guard is not exact manifold invariance; residual cross-block
leakage is small and measured numerically.)
\end{lemma}

This is what sets the precision-weighting regime: the reversed precision must be strong enough to move
the teacher along novel memory directions, yet weak enough that the self-pull wins normal to the
manifold.

\subsection{An exact free-energy saddle}
Both forces in Theorem~\ref{thm:prioritized} are gradients, so it is natural to ask when the coupled
$(\xT,\xS)$ flow is generated by a \emph{single} potential. It is, at one precise condition.

\begin{proposition}[Free-energy saddle]
\label{prop:saddle}
The sleep dynamics are generated by the single functional $\Phi=\FT-\FS$ --- the teacher descending
$\Phi$, the student ascending it --- if and only if
\begin{equation}
\label{eq:saddle}
\piST=-\piTS.
\end{equation}
Equivalently (Theorem~\ref{thm:prioritized}(ii)), the teacher ascends the student's surprise at exactly
the rate $\piTS$ that defines it: the ascent's coefficient $|\piST|/\piTS$ equals one. At this point the
pair plays a zero-sum game on one potential --- the student descending $\FS$
(Proposition~\ref{prop:descent}), the teacher climbing it --- a minimax whose shape is reminiscent of
the agent--environment saddle of active inference, though we do not claim identity of the two
\cite{friston2017process,millidge2021whence}.
\end{proposition}

\begin{remark}[Scope: subspace, not named episodes]
\label{rem:scope}
The substrate is linear, so a teacher's memories span a flat subspace and what transfers is that
subspace --- an effective rank, not a labelled list of patterns. Episodic, one-pattern-per-step replay
would require an added nonlinearity (e.g.\ a soft winner-take-all); we return to this in the Discussion.
\end{remark}

%============================================================================
\section{Results}
%============================================================================
We now show the mechanism at work and then that it composes. The first two subsections establish the
single-teacher primitive (rung~1) and its free-energy reading; the last two climb the ladder to the
union of subspaces (rung~2) and continual learning (rung~3). All simulations use the linear covariance
predictive-coding substrate of Section~\ref{sec:model}; parameters and diagnostics are in Materials and
methods.

\subsection{Prioritized transfer of unlearned directions}
\label{sec:res_transfer}
We start with the picture. Fig~\ref{fig:flow} takes a single teacher holding two memories and looks at
the teacher's state in the plane spanned by one direction the student has \emph{already} learned and one
it has \emph{not}. The reversed-precision flow tells the whole story at a glance: along the learned axis
the teacher is pinned --- the push there is null and the self-pull holds it in place --- while along the
novel axis it is driven outward, and a sample trajectory (red) rides the novel direction until the
student catches up. This is Theorem~\ref{thm:prioritized}(i) made visible: motion only along what the
student does not yet know. The figure is the two-network analogue of the attractor stream plots used for
single associative memories \cite{tang2023covpcn} and in our own earlier network
\cite{saighi2025autonomous}, and reading it the same way is intended.

\figplaceholder{Fig~\ref{fig:flow} source: a short script on the two-population model ---
project $\tauT\dot\xT$ from Eq.~\eqref{eq:teacher_reduced} onto the (learned, novel) plane; overlay one
transfer trajectory. Analogous to the CHN vector-field figure.}

\begin{figure}[h!]
\centering
\figplaceholder{\textbf{The reversed-precision flow pushes the teacher only along unlearned
directions.} Teacher state in a 2D slice: horizontal axis a memory direction the student has learned,
vertical axis one it has not. Arrows show $\tauT\dot\xT$ in replay ($\piST<0$). Along the learned axis
the flow is inward/pinned (push $=0$, self-pull holds); along the novel axis it is outward (push
$=|\piST|\,n\,c$). Red: one transfer trajectory, riding the novel axis until learning flattens it.}
\caption{\textbf{Prioritized replay is visible as a flow field.} See body text,
Section~\ref{sec:res_transfer}.}
\label{fig:flow}
\end{figure}

The quantitative claim is that transfer acquires the teacher's whole memory subspace, prioritizing the
directions the student lacks, and then stops. Fig~\ref{fig:staircase} shows it. The novelty spectrum ---
the eigenvalues of $\NS$ restricted to the teacher's memory manifold, $\UT\tr\NS\UT$ --- starts high on
every stored direction and falls to zero one direction at a time as the student learns them, producing a
descending ``staircase'' (left). The accompanying alignment heatmap (right) shows the teacher's state
progressively filling the memory manifold as the drive on each direction extinguishes. Crucially, the
number of steps is the memory set's \emph{effective rank}, not its pattern count: $P=6$ orthonormal
patterns give six steps, but $P$ correlated patterns of rank three give only three, because learning one
direction explains its correlated neighbours at once. This is the linear model transferring a
\emph{subspace} (Remark~\ref{rem:scope}). We report reliable completion of the manifold across the
regimes we tested; the exact step count and its ordering are a spectral, not a temporal, effect, and the
apparent sharpness of the steps is partly a plotting artefact of the timescale.

\begin{figure}[h!]
\centering
\figplaceholder{\textbf{Left:} restricted novelty spectrum $\mathrm{eig}(\UT\tr\NS\UT)$ vs.\ replay time:
a descending staircase, one step per learned direction; overlaid, orthonormal ($P{=}6\to 6$ steps) vs.\
correlated rank-3 ($\to 3$ steps). \textbf{Right:} heatmap of $\xT$--manifold alignment over time, filling
as the drive extinguishes. Source: \texttt{notebooks/two\_network/01\_single\_run.py},
\texttt{05\_findingH\_subspace.py}.}
\caption{\textbf{Transfer acquires the memory subspace and self-terminates; steps count effective rank.}
See Section~\ref{sec:res_transfer}.}
\label{fig:staircase}
\end{figure}

\subsection{The wake/sleep saddle picture}
\label{sec:res_saddle}
Why does the replay prioritize, and why does it stop? Section's theory answers both, and
Fig~\ref{fig:saddle} shows the answer directly. It plots cross-sections of the free-energy landscape the
teacher moves on. Along a direction the student can already predict, the landscape is a \emph{bowl}: the
teacher sits at the bottom, held on the manifold. Along a direction the student cannot predict, it is a
\emph{hill}: the reversed precision turns that axis into an unstable ascent, and the teacher slides up it
--- into the novel memory. As the student learns the direction, the hill relaxes back into a bowl, and
the drive is gone. Learning literally reshapes the landscape from hill to bowl, one direction at a time.

At the matching condition $\piST=-\piTS$ (Proposition~\ref{prop:saddle}) these cross-sections are two
faces of one potential $\Phi=\FT-\FS$: a genuine saddle, minimized by the teacher over learned
directions and maximized over novel ones. We verify the condition quantitatively by measuring the
circulation of the coupled flow, the obstruction to its being a single gradient; it vanishes precisely
at $\piST=-\piTS$ and grows linearly on either side, matching $|\piST+\piTS|\sqrt{d}$ to machine
precision.

Two further checks confirm the mechanism is the one we claim, and we relegate their figures to Supporting
information. First, stability: sweeping the strength $|\piST|$ shows the teacher consolidating
on-manifold below the guard of Lemma~\ref{lem:normal} and abruptly chasing noise above it, with the
turnover at the predicted value. Second, selection is spectral, not a race in time: across the entire
fast-student regime a known direction is explained away while a novel one is not, and the discrimination
collapses only as $\tauS\to\tauT$, i.e.\ only when the student is no longer fast --- so it is what the
student's weights can null, not raw speed, that decides.

\begin{figure}[h!]
\centering
\figplaceholder{Cross-sections of $\Phi=\FT-\FS$ at $\piST=-\piTS$: a bowl along a predictable direction
(teacher pinned) vs.\ a hill along a novel direction (teacher driven up it); a middle panel shows a hill
relaxing to a bowl as the direction is learned. Source:
\texttt{notebooks/two\_network/02\_findingD\_saddle.py}, \texttt{06\_eigenspace\_geometry.py}.}
\caption{\textbf{Replay is ascent on a hill that learning turns into a bowl.} At $\piST=-\piTS$ the two
faces belong to one saddle potential $\Phi=\FT-\FS$. See Section~\ref{sec:res_saddle}.}
\label{fig:saddle}
\end{figure}

\subsection{Merging memories by interleaving}
\label{sec:res_union}
Transfer is a primitive, and primitives compose. If a student is coupled to two teachers, each holding
its own memories, can it acquire the \emph{union} of their subspaces? It can, provided the two teachers
are replayed one at a time rather than together. We rehearse teacher $T_1$ for a bout, then $T_2$ for a
bout, and alternate; each bout is a transfer step of the kind analysed above, with only one interface
active. Over bouts the student accumulates $\mathcal{U}_1+\mathcal{U}_2$.

The subtlety, and the reason for interleaving, is \emph{crosstalk} between correlated memories.
Fig~\ref{fig:sawtooth} shows it as a sawtooth. Rehearsing $T_1$ drives the residual on $T_1$'s memory to
zero but, because the two memory sets overlap, nudges the residual on $T_2$'s memory up; rehearsing $T_2$
then nulls its own and nudges $T_1$'s back up. Left there, this is exactly the interference that makes
sequential learning forget. But the bumps shrink each round --- each rehearsal re-excites and re-cancels
the correlated perturbation the other one caused --- and the sawtooth decays to zero: both subspaces are
held at once. This is interleaved rehearsal, the classical cure for catastrophic forgetting, arising here
as the natural way to compose two transfer steps.

\begin{figure}[h!]
\centering
\figplaceholder{Union deficit and per-teacher residuals vs.\ rehearsal bout: a decaying sawtooth ---
rehearsing $T_1$ nulls its memory and bumps $T_2$'s, and vice versa, the bumps decaying to zero as both
subspaces are jointly held. Source:
\texttt{notebooks/subspace\_addition/01\_interleaved\_single\_run.py},
\texttt{02\_interleaved\_edge\_cases.py}.}
\caption{\textbf{Interleaved transfer merges two memory subspaces and cancels crosstalk.} The sawtooth
decays to zero: $\mathcal{U}_1+\mathcal{U}_2$ is retained. See Section~\ref{sec:res_union}.}
\label{fig:sawtooth}
\end{figure}

\subsection{Continual learning without catastrophic forgetting}
\label{sec:res_continual}
The top rung assembles transfer into a loop that absorbs a \emph{stream} of memories
(Algorithm~\ref{alg:continual}). Three networks divide the labour: a fast \textbf{buffer} that takes in
one new memory at a time (and is the only network that ever learns from external data), a transient
\textbf{synthesis} workspace, and a slow \textbf{storage} that accumulates the consolidated result. Each
cycle writes the new memory into the buffer, \emph{consolidates} buffer-and-storage into synthesis by
interleaved transfer (Section~\ref{sec:res_union}), then \emph{downloads} synthesis into storage. The
interleaving is what lets storage's existing memories be rehearsed alongside the new one, so old
correlated memories are re-excited and preserved rather than overwritten.

\begin{algorithm}[h]
\caption{Continual learning by prioritized transfer}
\begin{algorithmic}[1]
\State \textbf{Input:} stream of memories $m_1,m_2,\dots$; empty storage $\WS^{\mathrm{sto}}$
\For{each arriving memory $m_k$}
  \State \textbf{Write:} store $m_k$ in the fast buffer (one-shot), overwriting the previous
  \State \textbf{Consolidate:} interleaved transfer of \{buffer, storage\} $\to$ synthesis
         (alternate rehearsing buffer and storage; Section~\ref{sec:res_union})
  \State \textbf{Download:} transfer synthesis $\to$ storage
\EndFor
\State \textbf{Return:} storage holding $\mathrm{span}\{m_1,\dots,m_k\}$ up to capacity $d-1$
\end{algorithmic}
\label{alg:continual}
\end{algorithm}

Fig~\ref{fig:retention} shows the payoff. As memories stream in, we track a retention matrix: for each
stored memory (row), how well it is still recovered after each cycle (column). Under the loop, the matrix
fills and stays filled --- every memory acquired so far remains retrievable, and the stored subspace
grows one memory per cycle up to capacity. A buffer-only control, which simply writes each new memory
without the consolidate/download loop, retains only the most recent: its retention matrix is a single
bright diagonal against a forgotten background --- catastrophic forgetting in its purest form. The
difference is the loop, and the loop is nothing but transfer, applied twice per cycle.

\begin{figure}[h!]
\centering
\figplaceholder{\textbf{Left:} retention matrix under Algorithm~\ref{alg:continual} (memory $\times$
cycle) --- fills and stays filled; stored subspace grows one memory per cycle. \textbf{Right:} buffer-only
control --- only the latest memory retained (bright diagonal, forgotten elsewhere). Source:
\texttt{notebooks/continual\_learning/01\_continual\_single\_stream.py}.}
\caption{\textbf{The transfer loop retains a correlated stream where a buffer-only control forgets.}
See Section~\ref{sec:res_continual}.}
\label{fig:retention}
\end{figure}

%============================================================================
\section{Discussion}
%============================================================================
We have shown that a single signed precision turns a pair of predictive-coding memories from a recall
circuit into a replay circuit, and that in replay the downstream network's own prediction error steers
the upstream network into exactly the memories still missing downstream. The priority is not scheduled;
it is the gradient of the student's surprise (Theorem~\ref{thm:prioritized}), and at a precise
precision-matching condition the coupled dynamics are an exact free-energy saddle
(Proposition~\ref{prop:saddle}). Transfer is a primitive that composes into a union of subspaces and,
from there, into a continual-learning loop.

\paragraph{What is exact and what is only suggestive.} The variational free-energy reading is exact where
we use it: the student's inference and learning are gradient descent on $\FS$
(Proposition~\ref{prop:descent}), and the settled student's free energy is the novelty score the teacher
climbs (Corollary~\ref{cor:surprise}). We are deliberately more cautious with two larger claims. First,
that the loop minimizes a free energy \emph{integrated over the whole data stream} is a conjecture, not a
result --- we prove per-step descent, not a global objective. Second, the resemblance of the
$\piST=-\piTS$ saddle to the agent--environment minimax of active inference is structural, a matter of
shape; we do not claim the reversed precision implements an expected-free-energy computation
\cite{millidge2021whence}. Flagging these keeps the free-energy story honest.

\paragraph{Biological reading.} The architecture maps naturally onto complementary learning systems: a
fast teacher that already holds an experience (hippocampus) trains a slow, initially empty student
(neocortex) offline, and does so preferentially for what the neocortex has not yet consolidated
\cite{mcclelland1995cls,kumaran2016what,wilson1994reactivation,diekelmann2010memory}. The wake/sleep
distinction reduces to the sign of one interface precision, and prioritization requires no homunculus
choosing what to replay --- the error does the choosing, as in prioritized experience replay
\cite{schaul2016prioritized} but here derived from the network's own free energy. This complements our
earlier single-network account, where autonomous replay was produced instead by plastic self-inhibition
shallowing visited attractors \cite{saighi2025autonomous}: there one network recovered its own memories;
here two networks transfer them, and the priority signal is made explicit.

\paragraph{Limitations and outlook.} The price of analytic transparency is linearity, and with it the
scope of Remark~\ref{rem:scope}: transfer moves a memory \emph{subspace}, an effective rank, not a set of
named episodes. Episodic, one-pattern-per-bout replay --- the regime in which the biological analogy is
sharpest --- would need a nonlinearity such as a soft winner-take-all to select a single pattern per
replay, and building that on the same reversed-precision skeleton is the natural next step. Manifold
completion of the fully coupled dynamics is established empirically across the regimes we tested rather
than proven, and a stochastic-approximation treatment of the noise-seeded exploration remains open. None
of these affects the core claims, which are the exact ones above.

%============================================================================
\section{Conclusion}
%============================================================================
Prioritized memory transfer is a small, ownable primitive: reverse one interface precision, and a
downstream network's prediction error drives an upstream network to replay exactly what the downstream
network still lacks, stopping itself when the gap is closed. Because the mechanism is a gradient on the
downstream network's own surprise, priority comes for free; because the substrate is a linear
predictive-coding memory, the whole thing is exactly analyzable; and because transfer is a primitive, it
composes --- into a union of memory subspaces, and into a continual-learning loop that resists
catastrophic forgetting. The account is compact enough to prove and biologically suggestive enough to
motivate the nonlinear, episodic extensions it points to.

%============================================================================
\section*{Materials and methods}
%============================================================================
\paragraph{Substrate and faithfulness.} Each population is a linear covariance predictive-coding network
with one tied weight matrix $W$ (zero diagonal, enforced at construction and after every update):
top-down prediction uses $M=I-W$, bottom-up error projection uses $M\tr$, and the teacher--student
interface is identity in both directions. Memories are stored in the teacher as $\MT m_p=0$
($\max_p\|\MT m_p\|<10^{-4}$ for $P\le d-1$). The student's perception equals $-\nabla_{\xS}\FS$ and its
learning $-\nabla_{\WS}\FS$ (autograd-verified to $<10^{-8}$).

\paragraph{Proofs.} Proposition~\ref{prop:descent} (state and projected-weight gradient descent),
Lemma~\ref{lem:novelty} (fast-student elimination and $\NS$), Corollary~\ref{cor:surprise} (the surprise
identity, via per-direction summation $n^2+n(1-n)=n$ and, independently, the envelope theorem), and
Theorem~\ref{thm:prioritized} (the three readings of the push) are proved in full, every matrix step
justified, in Appendix~\ref{app:proofs}. Lemma~\ref{lem:normal} (the guard $|\piST|<\piT\sigma_{\min}^2$)
and Proposition~\ref{prop:saddle} (the saddle at $\piST=-\piTS$) are established with the numerical
diagnostics below.

\paragraph{Numerical diagnostics.} We verify: the surprise identity
$\FS(\xS^*,\xT)=\tfrac{\piTS}{2}\xT\tr\NS\xT$ (error $\sim10^{-15}$); the saddle circulation
$|\piST+\piTS|\sqrt{d}$, vanishing only at $\piST=-\piTS$; the off-manifold growth eigenvalue and terminal
manifold occupancy, both turning over at the guard; the restricted novelty spectrum
$\mathrm{eig}(\UT\tr\NS\UT)$ (the staircase); and, for the union and continual loops, per-bout residuals
and the retention matrix. Adiabatic (fast-student) elimination is used throughout for the reduced teacher
dynamics of Eq.~\eqref{eq:teacher_reduced}. Simulations use the timescale ordering
$\tauS\ll\tauT\ll 1/\eta$, isotropic teacher noise $\xi$, and a fixed teacher norm; full parameter tables
accompany each experiment notebook.

%============================================================================
\section*{Supporting information}
%============================================================================
\textbf{S1. Stability and timescale sweeps.} On-manifold consolidation vs.\ noise-chasing across
$|\piST|$ (the guard turnover); spectral, not temporal, selection across $\tauS/\tauT$.
\textbf{S2. Effective-rank detail.} Staircase step count vs.\ pattern correlation.
\textbf{S3. Dendritic (stop-gradient) variant.} The mechanism under Tang et al.'s bio-plausible
stop-gradient relaxation \cite{tang2023covpcn}: transfer still completes, at the cost of a
$\sim\!2\times$ looser residual floor and a non-symmetric novelty operator.
\textbf{S4. Noise-free slow-mixing variant.} Transfer driven by integrator imprecision rather than
injected noise.

%============================================================================
\appendix
\section{Proofs}
\label{app:proofs}
%============================================================================
We give the full proofs of the four analytic results stated in the body. Every matrix step is
justified explicitly. We use the house identities $M=I-W$, $S=M\tr M$, $\epsT=\MT\xT$, $\epsS=\MS\xS$,
$\epsTS=\xS-\xT$ (Eqs.~\ref{eq:mismatch}--\ref{eq:interface}), and the free energies
$\FS=\tfrac{\piTS}{2}\|\epsTS\|^2+\tfrac{\piS}{2}\|\epsS\|^2$, $\FT=\tfrac{\piT}{2}\|\epsT\|^2$
(Eq.~\ref{eq:freeenergy}). Two algebra facts recur: \textbf{(R1)} $(AB)\tr=B\tr A\tr$; and
\textbf{(R2)} for $g=\tfrac12\|u\|^2$, a perturbation $u\to u+\delta$ gives $g\to g+u\tr\delta+O(\|\delta\|^2)$,
so $\nabla g=u$.

\paragraph{Timescales.} Two distinct ``student fixed'' assumptions are used, on the nested separation
$\tauS\ll\tauT\ll1/\eta$. Lemma~\ref{lem:novelty} uses the student \emph{state} being \emph{fast}: $\xS$
is already settled at every teacher configuration ($\dot\xS=0$). Theorem~\ref{thm:prioritized} uses the
student \emph{weights} being \emph{slow}: over the teacher's exploration $\WS$, and hence $\SS$ and
$\NS$, are frozen matrices. The two are not the same assumption.

\subsection{Proposition~\ref{prop:descent}: inference and learning descend \texorpdfstring{$\FS$}{FS}}
\label{app:prop1}

\paragraph{Part A --- inference.} We first derive the one differentiation rule. For fixed $A$ and
$g(x)=\tfrac12\|Ax\|^2=\tfrac12(Ax)\tr(Ax)$, perturb $x\to x+\delta$ and expand exactly:
\[
g(x+\delta)=\tfrac12(Ax)\tr(Ax)+(Ax)\tr(A\delta)+\tfrac12(A\delta)\tr(A\delta).
\]
The linear term is $(Ax)\tr A\delta=x\tr A\tr A\,\delta$ by (R1); its coefficient transposed gives
\begin{equation}
\label{eq:gradnorm}
\nabla_x\,\tfrac12\|Ax\|^2=A\tr A\,x.
\end{equation}
(The $\tfrac12$ cancels the factor $2$ from the two equal cross terms.) Apply
Eq.~\eqref{eq:gradnorm} to each term of $\FS$, holding $\xT$ fixed. The interface term
$\tfrac{\piTS}{2}\|\xS-\xT\|^2$ is the form with $A=I$ and a constant shift (which drops), giving
$\piTS\,\epsTS$. The self term $\tfrac{\piS}{2}\|\MS\xS\|^2$ is the form with $A=\MS$, giving
$\piS\MS\tr\epsS$. Summing,
\[
\nabla_{\xS}\FS=\piTS\,\epsTS+\piS\,\MS\tr\epsS,
\]
and the student's state equation \eqref{eq:student} is exactly $\tauS\dot\xS=-\nabla_{\xS}\FS$: inference
is exact gradient descent on $\FS$.

\paragraph{Part B --- learning.} Only the self term $\tfrac{\piS}{2}\|\xS-\WS\xS\|^2$ depends on $\WS$.
Writing it as $\tfrac{\piS}{2}\sum_i\varepsilon_{S,i}^2$ with
$\varepsilon_{S,i}=x_{S,i}-\sum_k(\WS)_{ik}x_{S,k}$, and using
$\partial(\WS)_{ik}/\partial(\WS)_{ab}=\delta_{ia}\delta_{kb}$,
\[
\frac{\partial\varepsilon_{S,i}}{\partial(\WS)_{ab}}=-\delta_{ia}\,x_{S,b}
\quad\Longrightarrow\quad
\frac{\partial\FS}{\partial(\WS)_{ab}}=\piS\sum_i\varepsilon_{S,i}(-\delta_{ia}x_{S,b})=-\piS\,\varepsilon_{S,a}\,x_{S,b},
\]
the $\delta_{ia}$ collapsing the sum. Assembling the entries into a matrix gives the outer product
$\nabla_{\WS}\FS=-\piS\,\epsS\,\xS\tr$. The zero-diagonal constraint restricts admissible updates to the
subspace $Z$ of zero-diagonal matrices; the orthogonal projection onto $Z$ is $P_0(X)=X-\mathrm{diag}(X)$,
which is idempotent ($P_0^2=P_0$) and self-adjoint ($\langle P_0X,Y\rangle=\langle X,P_0Y\rangle$) for the
Frobenius inner product $\langle X,Y\rangle=\mathrm{tr}(X\tr Y)$. The learning rule follows the projected
gradient $\dot\WS=-\eta\,P_0(\nabla_{\WS}\FS)$, so with $G=\nabla_{\WS}\FS$,
\[
\dot\FS=\langle G,\dot\WS\rangle=-\eta\,\langle G,P_0G\rangle=-\eta\,\langle P_0G,P_0G\rangle=-\eta\,\|P_0G\|_F^2\le0,
\]
inserting $P_0G=P_0(P_0G)$ then moving one $P_0$ across. Learning is projected gradient descent on $\FS$.
\hfill$\blacksquare$

\subsection{Lemma~\ref{lem:novelty}: the novelty operator}
\label{app:lem1}

\paragraph{Step 1--2 --- solve the fast-student steady state.} ``Fast student'' means
$\nabla_{\xS}\FS=0$, i.e.\ $\piTS(\xS-\xT)+\piS\SS\xS=0$. Collecting $\xS$ on the left with
$B:=\piTS I+\piS\SS$,
\[
B\,\xS=\piTS\,\xT
\quad\Longrightarrow\quad
\xS^{*}=\piTS\,B^{-1}\xT.
\]

\paragraph{Step 3 --- $B$ is invertible.} If $\SS v=\mu v$ then $Bv=(\piTS+\piS\mu)v$, so $B$ shares the
eigenbasis of $\SS$ with eigenvalues $\piTS+\piS\mu\ge\piTS>0$ (using $\mu\ge0$ from Step~5). None is
zero, so $B^{-1}$ exists and $B^{-1}v=(\piTS+\piS\mu)^{-1}v$.

\paragraph{Step 4 --- the leftover interface error.} Writing $\xT=B^{-1}B\xT$ to share a left factor,
\[
\epsTS=\xS^{*}-\xT=B^{-1}\big[\piTS I-(\piTS I+\piS\SS)\big]\xT=B^{-1}(-\piS\SS)\xT=-\piS\,B^{-1}\SS\,\xT.
\]

\paragraph{Step 5 --- $\SS$ is symmetric PSD, and learned $\Rightarrow$ zero.} For any $v$,
$v\tr\SS v=v\tr\MS\tr\MS v=\|\MS v\|^2\ge0$, so $\SS$ is PSD ($\mu\ge0$) and symmetric ($(\MS\tr\MS)\tr=\MS\tr\MS$).
A learned direction $u$ has $\MS u=0$, hence $\SS u=\MS\tr\MS u=0$.

\paragraph{Step 6--7 --- $\SS$ commutes with $B^{-1}$; the operator $\NS$.} Since $\SS
B=\piTS\SS+\piS\SS^2=B\SS$, sandwiching by $B^{-1}$ gives $B^{-1}\SS=\SS B^{-1}$. Sliding $\SS$ left in
Step~4,
\[
\epsTS=-\piS\,\SS\,B^{-1}\xT=-\NS\,\xT,\qquad
\NS=\piS\,\SS\,(\piTS I+\piS\SS)^{-1}.
\]
$\NS$ is symmetric (product of the commuting symmetric $\SS$ and $B^{-1}$: $(\SS B^{-1})\tr=B^{-1}\SS=\SS
B^{-1}$) and PSD (eigenvalues $n(\mu)=\piS\mu/(\piTS+\piS\mu)\ge0$).

\paragraph{Step 8 --- $\NS$ annihilates learned directions.} For learned $u$, Step~5 gives $\SS u=0$, so
$\NS u=\piS B^{-1}\SS u=0$ and $\epsTS=-\NS u=0$: a teacher along a learned direction produces no drive.
\hfill$\blacksquare$

\paragraph{Complement identity.} From $B B^{-1}=I$, i.e.\ $\piTS B^{-1}+\piS\SS B^{-1}=I$, and the
definition of $\NS$, we get $\piTS B^{-1}=I-\NS$, hence
\[
\xS^{*}=(I-\NS)\,\xT,\qquad \epsTS=(I-\NS)\xT-\xT=-\NS\xT.
\]
The student copies the predicted part ($I$) and leaves the novel part ($-\NS$): per direction it keeps
the fraction $1-n_k$ and leaves $n_k$ behind as error.

\subsection{Corollary~\ref{cor:surprise}: the surprise identity}
\label{app:cor1}

Let the student settle ($\xS=\xS^{*}$) with the teacher at $\xT$, and read off
$\FSeq(\xT):=\FS(\xS^{*}(\xT),\xT)$; since $\FS$ is a strictly convex quadratic in $\xS$ (Hessian $B\succ0$),
this equals $\min_{\xS}\FS$.

\paragraph{Direct computation.} Plug in $\xS^{*}=(I-\NS)\xT$ and $\epsTS=-\NS\xT$. The interface term,
using $\NS$ symmetric,
\[
\tfrac{\piTS}{2}\|\epsTS\|^2=\tfrac{\piTS}{2}\,\xT\tr\NS^2\,\xT.
\]
The self term, using $\|\epsS\|^2=(\xS^{*})\tr\SS\,\xS^{*}$ and the key move
$\piS\SS(I-\NS)=\piTS(\piS\SS B^{-1})=\piTS\NS$ (complement identity),
\[
\tfrac{\piS}{2}\|\epsS\|^2=\tfrac{\piS}{2}\,\xT\tr(I-\NS)\SS(I-\NS)\xT=\tfrac{\piTS}{2}\,\xT\tr(\NS-\NS^2)\xT.
\]
Adding, the $\NS^2$ terms cancel:
\[
\FSeq(\xT)=\tfrac{\piTS}{2}\,\xT\tr\big[\NS^2+\NS-\NS^2\big]\xT=\tfrac{\piTS}{2}\,\xT\tr\NS\,\xT.
\]
Per direction $\xT=\sum_k c_k u_k$: interface costs $n_k^2$, self costs $n_k(1-n_k)$, summing to $n_k$;
so $\FSeq=\tfrac{\piTS}{2}\sum_k n_k c_k^2$ --- novelty priced on where the state sits.
\hfill$\blacksquare$

\paragraph{Gradient (envelope route).} For $g(\theta)=\min_x f(x,\theta)=f(x^{*}(\theta),\theta)$ the chain
rule gives $g'(\theta)=\underbrace{(\partial f/\partial x)|_{x^{*}}}_{=0}\,dx^{*}/d\theta+(\partial
f/\partial\theta)|_{x^{*}}$, the first factor vanishing by the definition of the minimizer. Here
$\theta=\xT$ and only the interface term carries $\xT$ explicitly, with $\xT$-gradient $\piTS(\xT-\xS)$;
freezing $\xS=\xS^{*}$,
\[
\nabla_{\xT}\FSeq=\piTS(\xT-\xS^{*})=-\piTS\,\epsTS=\piTS\,\NS\,\xT,
\]
matching the direct computation.

\subsection{Theorem~\ref{thm:prioritized}: prioritized transfer}
\label{app:thm1}

\paragraph{Step 1 --- read off the push.} The teacher in sleep obeys
$\tauT\dot\xT=-\piT\MT\tr\epsT+\piST\epsTS+\xi$. Substituting $\MT\tr\epsT=\ST\xT$, the fast-student
$\epsTS=-\NS\xT$ (Lemma~\ref{lem:novelty}), dropping $\xi$, and writing $\piST=-|\piST|$,
\[
\tauT\dot\xT=\underbrace{-\piT\ST\xT}_{\text{self-pull}}+\underbrace{|\piST|\NS\xT}_{u},
\qquad u=|\piST|\,\NS\,\xT.
\]

\paragraph{Step 2 --- per-direction amplification.} $\NS$ is symmetric, so $\NS u_k=n_k u_k$ on an
orthonormal eigenbasis, $n_k\in[0,1)$, $n_k=0$ iff learned. With $c_k=u_k\tr\xT$ and under the push
alone ($\dot\xT=u/\tauT$),
\[
\dot c_k=u_k\tr\dot\xT=\frac{|\piST|}{\tauT}\,u_k\tr\NS\xT=\frac{|\piST|}{\tauT}\,n_k\,c_k,
\]
using $u_k\tr\NS=n_k u_k\tr$. Learned directions ($n_k=0$) are frozen; novel ones grow, faster the more
novel.

\paragraph{Step 3 --- steepest ascent of the student's surprise.} By Corollary~\ref{cor:surprise},
$\FSeq(\xT)=\tfrac{\piTS}{2}\xT\tr\NS\xT=\tfrac{\piTS}{2}\sum_k n_k c_k^2$, whose gradient (by perturbation,
as in Eq.~\ref{eq:gradnorm}, $\NS$ symmetric) is $\nabla_{\xT}\FSeq=\piTS\NS\xT$. Hence
\[
u=|\piST|\,\NS\,\xT=\frac{|\piST|}{\piTS}\,\nabla_{\xT}\FSeq,
\]
a positive multiple of the surprise gradient. By Cauchy--Schwarz $\nabla\FSeq\cdot e\le\|\nabla\FSeq\|$
with equality iff $e\parallel\nabla\FSeq$: of all directions, the push raises the student's surprise
fastest (a local steepest ascent).

\paragraph{Step 4 --- climbing and self-limiting.} Along the push,
\[
\frac{d\FSeq}{dt}=\nabla\FSeq\cdot\dot\xT=\frac{\piTS|\piST|}{\tauT}\,\|\NS\xT\|^2\ge0,
\]
zero only when $\NS\xT=0$. And $n(\mu)\le(\piS/\piTS)\mu$ gives
$|\piST|n_k\le(|\piST|\piS/\piTS)\|\MS u_k\|^2$, so as the student learns a direction
($\|\MS u_k\|\to0$, Proposition~\ref{prop:descent}) the push there vanishes: the steering self-terminates
direction by direction.

\paragraph{Step 5 --- the saddle condition (preview of Proposition~\ref{prop:saddle}).} The self-pull is
$-\nabla_{\xT}\FT$ with $\nabla_{\xT}\FT=\piT\ST\xT$, and the push is $(|\piST|/\piTS)\nabla_{\xT}\FSeq$, so
\[
\tauT\dot\xT=-\nabla_{\xT}\FT+\frac{|\piST|}{\piTS}\,\nabla_{\xT}\FSeq,
\]
which is $-\nabla\Phi$ for the single potential $\Phi=\FT-\FSeq$ exactly when the ascent coefficient is
one: $|\piST|=\piTS$, i.e.\ $\piST=-\piTS$. The ascent is local, and the push rescales but does not seed
an unoccupied direction ($c_k=0\Rightarrow\dot c_k=0$); noise supplies that seed in the full model.
\hfill$\blacksquare$

%============================================================================
\bibliography{references}
%============================================================================

\end{document}
